% English appears first; Thai follows directly below. All original marks are retained.
\bpart{1}{(25 คะแนน) ข้อกำหนดและความมั่นคงปลอดภัย}{(25 marks) Requirements and Security}
\bq{1}{วิเคราะห์ความต้องการของผู้มีส่วนได้ส่วนเสีย}{Analyse stakeholder requirements}{8}
\begin{prompttable}
\prow{มหาวิทยาลัยกำลังเตรียมระบบลงทะเบียน workshop ผู้ประสานงานส่งข้อความความต้องการสี่ข้อให้ทีมพัฒนา แต่ทีมยังไม่ได้ยืนยันรายละเอียดอื่นนอกเหนือจากข้อความที่ให้ ให้วิเคราะห์แต่ละข้อแยกกัน โดยใช้หัวข้อต่อไปนี้}
{A university is preparing a workshop registration system. Its coordinator gives the development team four requirement statements. No additional details have yet been agreed beyond those stated. Analyse each statement independently using the following criteria.}
\prow{จำแนกประเภทเป็น functional requirement, non-functional requirement หรือ business/domain rule พร้อมเหตุผล อธิบายจุดประสงค์ ประเมินคุณภาพและจุดที่ยังไม่ชัด เสนอวิธีทดสอบหรือวัดพร้อมผลที่คาดหวัง แล้วปรับข้อความหรืออธิบายเหตุผลที่ควรคงไว้ ตัวเลขหรือเงื่อนไขที่เสนอเพิ่มต้องระบุว่ารอยืนยันกับผู้เกี่ยวข้อง}
{Classify it as a functional requirement, non-functional requirement, or business/domain rule and justify the classification. Explain its purpose, assess its quality and ambiguities, propose a test or measurement with the expected result, then revise the statement or justify retaining it. Identify any additional conditions or numerical targets as proposals requiring stakeholder confirmation.}
\end{prompttable}
\bsub{1.1}{ดู workshop ที่เปิดรับลงทะเบียน}{View workshops open for registration}{2}
\pair{“ระบบต้องให้นักศึกษาที่เข้าสู่ระบบแล้วดู workshop ทั้งหมดที่เปิดรับลงทะเบียนได้” วิเคราะห์ข้อความนี้ตามหัวข้อของข้อ 1 ให้ครบ}
{“The system shall let a signed-in student view all workshops that are open for registration.” Analyse this statement using all the criteria in Question 1.}
\answerlabel\answerfill
\newpage
\continued{ข้อ 1 (ต่อ): ใช้หัวข้อวิเคราะห์เดียวกับหน้าแรกของข้อ 1}{Question 1 (continued): Apply the same analysis criteria}
\bsub{1.2}{ความเร็วในการลงทะเบียน}{Registration speed}{2}
\pair{“กระบวนการลงทะเบียนควรรวดเร็ว” อธิบายว่าข้อความนี้บอกอะไรและยังขาดข้อมูลใด เสนอวิธีตรวจและผลที่คาดหวังซึ่งสอดคล้องกับการตีความของคุณ หากเพิ่มเกณฑ์ขึ้นเอง ให้ระบุว่าเป็นข้อเสนอที่ยังต้องยืนยัน}
{“The registration process should be quick.” Explain what this statement expresses and what information is missing. Propose a check and expected result consistent with your interpretation. Mark any criteria you introduce as recommendations that still require confirmation.}
\answerlabel\blankbox{6.6cm}
\bsub{1.3}{ข้อจำกัดตามนโยบายมหาวิทยาลัย}{University policy limit}{2}
\pair{“ตามนโยบายมหาวิทยาลัย นักศึกษาแต่ละคนลงทะเบียน workshop ได้ไม่เกิน 3 รายการต่อภาคการศึกษา” วิเคราะห์ตามหัวข้อเดิม โดยไม่เปลี่ยนนโยบายที่ให้ และระบุข้อมูลที่ควรถามเพิ่มก่อนนำไปพัฒนาหรือทดสอบ}
{“University policy limits each student to three workshop registrations per semester.” Apply the same analysis criteria without changing the stated policy. Identify any details that should be clarified before implementation or testing.}
\answerlabel\answerfill
\newpage
\continued{ข้อ 1 (ต่อ): วิเคราะห์ข้อความข้อสุดท้าย}{Question 1 (continued): Analyse the final statement}
\bsub{1.4}{ยกเลิกการลงทะเบียนและดูรายชื่อ}{Cancel registrations and view participant lists}{2}
\pair{“ระบบต้องให้นักศึกษายกเลิกการลงทะเบียน และให้ผู้ประสานงานดูรายชื่อผู้เข้าร่วมได้” วิเคราะห์ประเภท จุดประสงค์ คุณภาพ วิธีตรวจ และสิ่งที่ควรปรับ โดยให้เหตุผลจากข้อความที่ได้รับ หากแยกเป็นหลายข้อ ต้องรักษาความต้องการเดิมให้ครบ และไม่เพิ่มสิทธิ์หรือเงื่อนไขใหม่โดยถือว่าเป็นข้อเท็จจริง}
{“The system shall let students cancel a registration and let coordinators see participant lists.” Analyse its type, purpose, quality, verification and possible revision, using evidence from the statement. If you split it into multiple requirements, preserve all the original intent. Do not treat newly introduced permissions or conditions as given facts.}
\answerlabel\answerfill
\newpage
\bq{2}{ประเมินข้อเสนอโครงการ WorkshopHub}{Evaluate the WorkshopHub proposal}{7}
\begin{prompttable}
\prow{ปัจจุบันนักศึกษาสมัคร workshop ของมหาวิทยาลัยผ่านแบบฟอร์มออนไลน์ที่แยกจากกัน ผู้ประสานงานรวบรวมคำตอบด้วยมือ บางครั้งจึงเกิดการลงทะเบียนซ้ำและยืนยันล่าช้า}
{Students currently register for campus workshops through separate online forms. Coordinators combine the responses manually, sometimes causing duplicate registrations and delayed confirmations.}
\prow{ทีมเสนอ WorkshopHub ซึ่งเป็นเว็บไซต์สำหรับนักศึกษาและผู้ประสานงาน นักศึกษาค้นหา workshop และลงทะเบียนได้ ส่วนผู้ประสานงานจัดการรายชื่อผู้เข้าร่วม ทีมคาดว่าจะลดข้อผิดพลาดและงานจัดการ รุ่นแรกครอบคลุม workshop ฟรีเท่านั้น การชำระเงินและใบรับรองจะทำภายหลัง ทีมวางแผนให้เสร็จก่อนภาคการศึกษาถัดไป}
{The team proposes WorkshopHub, a website for students and workshop coordinators. Students can browse and register for workshops, while coordinators manage participant lists. The team expects fewer errors and less administrative work. The first release supports free workshops only; payment and attendance certificates will come later. The team plans to finish before the next semester.}
\end{prompttable}
\bsub{2.1}{ประเมินองค์ประกอบของข้อเสนอ}{Evaluate the proposal elements}{4}
\pair{ประเมินข้อเสนอตามที่เขียนจริงให้ครบ 8 องค์ประกอบ ระบุว่าแต่ละส่วนมีแล้ว มีเพียงบางส่วน หรือยังไม่มี พร้อมยกข้อความหรือเหตุผลประกอบ อย่านำข้อเสนอที่คุณคิดเพิ่มไปนับว่ามีอยู่ในต้นฉบับ}
{Assess all eight elements in the proposal as written. Mark each as present, incomplete, or missing, and support your judgment with a phrase or reason from the proposal. Do not count your own recommendations as information already provided.}
\begingroup\renewcommand{\tabularxcolumn}[1]{m{#1}}
\begin{tabularx}{\linewidth}{|>{\raggedright\arraybackslash}m{4.1cm}|P|}\hline
\bilingualcell{Element}{องค์ประกอบ} & \bilingualcell{Status and evidence}{สถานะและหลักฐาน}\\\hline
Problem & \rule{0pt}{.95cm}\\\hline
Users & \rule{0pt}{.95cm}\\\hline
Solution & \rule{0pt}{.95cm}\\\hline
Expected value & \rule{0pt}{.95cm}\\\hline
Scope & \rule{0pt}{.95cm}\\\hline
Risk & \rule{0pt}{.95cm}\\\hline
Effort & \rule{0pt}{.95cm}\\\hline
Evidence of completion & \rule{0pt}{.95cm}\\\hline
\end{tabularx}\endgroup
\newpage
\continued{ข้อ 2 (ต่อ): ใช้ข้อเสนอ WorkshopHub จากหน้าก่อน}{Question 2 (continued): Use the preceding WorkshopHub proposal}
\bsub{2.2}{เขียนข้อเสนอฉบับปรับปรุง}{Write an improved proposal}{3}
\pair{เขียนข้อเสนอที่ช่วยให้ทีมตัดสินใจได้ว่าจะสร้างอะไร และช่วยให้ผู้ตรวจรับทราบว่างานเสร็จเมื่อใด คงปัญหา ผู้ใช้ และขอบเขตรุ่นแรกตามโจทย์ เติมหรือเสนอวิธียืนยันส่วนที่ยังขาด แยกข้อมูลเดิมออกจากคำแนะนำของคุณให้ชัด โดยไม่แต่งตัวเลขประมาณการขึ้นเป็นข้อเท็จจริง ตอบเป็นย่อหน้าต่อเนื่องหรือหัวข้อที่เชื่อมโยงกันได้}
{Revise the proposal so that a team can decide what to build and a reviewer can determine when it is complete. Preserve the original problem, users and first-release scope. Fill the gaps or propose how to establish the missing information. Distinguish the original facts from your recommendations; do not present invented estimates as agreed facts. You may write connected paragraphs or clearly related points.}
\answerlabel\answerfill
\newpage
\bq{3}{สิทธิ์การใช้งานตามบทบาท}{Role-Based Access Control}{4}
\begin{prompttable}
\prow{WorkshopHub กำหนดหนึ่ง role ต่อหนึ่งบัญชี ไม่มีการสืบทอดสิทธิ์ระหว่าง roles และปฏิเสธ operation ที่นโยบายไม่ได้อนุญาต Student ดู workshop ที่เผยแพร่และส่งคำขอลงทะเบียนของตนเองได้}
{WorkshopHub assigns exactly one role to each account. Roles do not inherit permissions, and any operation not granted by the policy is denied. A Student may view published workshops and submit their own registrations.}
\prow{Coordinator ดู workshop ที่เผยแพร่ สร้าง workshop แก้รายละเอียด workshop และส่งออกรายชื่อผู้เข้าร่วมได้ โจทย์ไม่ได้กำหนดข้อจำกัดตามเจ้าของ workshop เพิ่มเติม}
{A Coordinator may view published workshops, create workshops, edit workshop details, and export participant lists. No additional workshop-ownership restriction is specified.}
\end{prompttable}
\bsub{3.1}{เติมตารางสิทธิ์ให้ครบทุกช่อง}{Complete every permission cell}{2.5}
\pair{ใส่ Allow หรือ Deny ตามนโยบายข้างต้น ไม่สมมติว่าบทบาทหนึ่งได้รับสิทธิ์ของอีกบทบาทโดยอัตโนมัติ}
{Enter Allow or Deny according to the stated policy. Do not assume that one role automatically receives the permissions of another.}
\begin{tabularx}{\linewidth}{|P|M{2.5cm}|M{2.5cm}|}\hline
Operation & Student & Coordinator\\\hline
\bilingualcell{View published workshops}{ดู workshop ที่เผยแพร่} &\rule{0pt}{.85cm}&\\\hline
\bilingualcell{Submit own registration}{ส่งคำขอลงทะเบียนของตน} &\rule{0pt}{.85cm}&\\\hline
\bilingualcell{Create a workshop}{สร้าง workshop} &\rule{0pt}{.85cm}&\\\hline
\bilingualcell{Edit workshop details}{แก้ workshop} &\rule{0pt}{.85cm}&\\\hline
\bilingualcell{Export participant lists}{ส่งออกรายชื่อ} &\rule{0pt}{.85cm}&\\\hline
\end{tabularx}
\bsub{3.2}{การซ่อนปุ่มเพียงพอหรือไม่}{Is hiding the button sufficient?}{1.5}
\pair{ทีมซ่อนปุ่มส่งออกรายชื่อจาก Student แล้ว การทำเช่นนี้เพียงพอหรือไม่? อธิบายสิ่งที่ backend ต้องตรวจ และเสนอ direct-request test หนึ่งกรณีพร้อมผลที่คาดหวัง}
{The team hides the export button from Students. Is this sufficient? Explain what the backend must check, and propose one direct-request test with its expected result.}
\answerlabel\answerfill
\newpage
\bq{4}{วิเคราะห์ช่องโหว่ระบบส่งรายงาน}{Analyse a report-upload vulnerability}{6}
\begin{prompttable}
\prow{ระบบส่งงานของมหาวิทยาลัยให้นักศึกษาอัปโหลดรายงาน PDF แต่ตรวจชนิดไฟล์เฉพาะใน browser นักศึกษาคนหนึ่งข้ามการตรวจดังกล่าว แล้วส่ง script ที่ปลอมเป็น PDF เข้าสู่ server}
{A university assignment system allows students to upload PDF reports but checks the file type only in the browser. A student bypasses this check and sends a script disguised as a PDF to the server.}
\prow{Server เก็บไฟล์นั้นไว้ในโฟลเดอร์ที่อนุญาตให้รัน script เมื่อมีผู้เปิด URL ของไฟล์ server จึงรัน script และ script อ่านรายงานของนักศึกษาคนอื่นได้ ใช้เหตุการณ์ที่ระบุนี้ในการตอบ และแยกผลที่เกิดแล้วออกจากผลที่เพียงอาจเกิดขึ้น}
{The server stores the file in a folder that permits script execution. Opening its URL causes the server to execute the script, which reads other students' reports. Base your answer on this incident. Distinguish effects that have occurred from those that might merely be possible.}
\end{prompttable}
\bsub{4.1}{วิเคราะห์เหตุการณ์และขอบเขตความเชื่อถือ}{Analyse the incident and trust boundaries}{3}
\pair{ระบุผู้โจมตี สิ่งที่ต้องปกป้อง (assets) จุดที่ผู้โจมตีเข้าถึงระบบ (attack surface) และ trust boundary ที่ถูกข้าม อธิบาย security objective ที่เสียหายพร้อมเชื่อมกับข้อเท็จจริงในเหตุการณ์ อย่าตอบเพียงชื่อคำศัพท์โดยไม่มีเหตุผล}
{Identify the attacker, protected assets, attack surface and crossed trust boundaries. Explain which security objective is harmed and connect your reasoning to the incident. Naming terms without explaining their relevance is not sufficient.}
\answerlabel\answerfill
\newpage
\continued{ข้อ 4 (ต่อ): ใช้เหตุการณ์อัปโหลดจากหน้าก่อน}{Question 4 (continued): Use the preceding upload incident}
\bsub{4.2}{มาตรการป้องกันและการทดสอบเชิงลบ}{Controls and negative tests}{3}
\pair{เสนอ control และ negative test จำนวน 2 คู่ แต่ละคู่ต้องระบุว่าจะป้องกันปัญหาตรงใด วิธีทดสอบที่พยายามข้ามการป้องกัน และผลที่ควรได้ เลือกมาตรการที่แก้ความเสี่ยงในเหตุการณ์จริง ไม่ใช้การตรวจใน browser เพียงอย่างเดียวเป็นมาตรการทั้งหมด}
{Propose two control-and-negative-test pairs. For each pair, identify the weakness it addresses, describe a test that attempts to bypass the protection, and state the expected result. Choose controls relevant to the incident; browser-only validation must not be your entire defense.}
\pair{\textbf{คู่ที่ 1}}{\textbf{Pair 1}}\blankbox{7.7cm}
\pair{\textbf{คู่ที่ 2}}{\textbf{Pair 2}}\answerfill
\newpage
\bpart{2}{(25 คะแนน) การวางแผนและดำเนินโครงการ}{(25 marks) Project Planning and Execution}
\bq{5}{วางแผนโครงการ DormFix}{Plan the DormFix project}{10}
\pair{คุณเป็นผู้จัดการโครงการ DormFix ซึ่งใช้ Waterfall เพื่อสร้างเว็บแจ้งซ่อมหอพัก ทีมมี Aom (BA), Bank (UX/UI), Chai (Developer), Dao (Developer/Tester) ผู้จัดการหอพักเป็น Sponsor และหัวหน้าช่างเป็นตัวแทน User งานในตารางไม่ได้เรียงตามลำดับ ให้ใช้ข้อมูลนี้ทำข้อ 5.1-5.3}
{You are the project manager of DormFix, a Waterfall web application for dormitory repairs. Your team is Aom (BA), Bank (UX/UI), Chai (Developer) and Dao (Developer/Tester). The dormitory manager is the Sponsor, and the maintenance head represents the User. Tasks below are not in execution order. Use them for Questions 5.1-5.3.}
\begingroup\fontsize{9.5}{12.5}\selectfont
\renewcommand{\tabularxcolumn}[1]{m{#1}}
\begin{tabularx}{\linewidth}{|M{.55cm}|P|M{1.1cm}|M{.75cm}|}\hline
ID & \bilingualcell{Task}{งาน} & Owner & Days\\\hline
A & \bilingualcell{Code the request form from the approved designs}{เขียนแบบฟอร์มแจ้งซ่อมจากแบบที่อนุมัติ} & Chai & 4\\\hline
B & \bilingualcell{Submit the requirements for approval by the dormitory manager}{ส่ง requirements ให้ผู้จัดการหอพักอนุมัติ} & PM & 1\\\hline
C & \bilingualcell{Design the screens from the approved requirements}{ออกแบบหน้าจอจาก requirements ที่อนุมัติ} & Bank & 3\\\hline
D & \bilingualcell{Test the two coded parts together}{ทดสอบโปรแกรมสองส่วนร่วมกัน} & Dao & 2\\\hline
E & \bilingualcell{Write the requirements; this is the first task}{เขียน requirements เป็นงานแรก} & Aom & 3\\\hline
F & \bilingualcell{Submit both designs for approval by the maintenance head}{ส่งแบบทั้งสองให้หัวหน้าช่างอนุมัติ} & PM & 1\\\hline
G & \bilingualcell{Design the database from the approved requirements}{ออกแบบฐานข้อมูลจาก requirements ที่อนุมัติ} & Chai & 2\\\hline
H & \bilingualcell{Submit the tested system for acceptance by the dormitory manager}{ส่งระบบที่ทดสอบแล้วให้ผู้จัดการหอพักตรวจรับ} & PM & 1\\\hline
I & \bilingualcell{Code the technician dashboard from the approved designs}{เขียน dashboard ช่างจากแบบที่อนุมัติ} & Dao & 3\\\hline
\end{tabularx}\endgroup
\pair{ใช้ชื่อผู้ทำและจำนวนวันตามตารางในการวางแผน รายละเอียดการอนุมัติและคำว่า “จากแบบที่อนุมัติ” หรือ “ระบบที่ทดสอบแล้ว” เป็นข้อมูลที่ต้องใช้กำหนดความสัมพันธ์ของงาน ไม่จำเป็นต้องใช้ทุกช่องของแบบฟอร์มคำตอบ}
{Use the stated owner and duration for scheduling. Approval steps and phrases such as “from the approved designs” and “the tested system” determine task relationships. You do not need to fill every available answer box.}
\newpage
\continued{ข้อ 5 (ต่อ): ใช้รายการงาน A-I}{Question 5 (continued): Use Tasks A-I}
\bsub{5.1}{จัด WBS ตาม milestone ที่บรรลุ}{Organize a milestone-based WBS}{3}
\begin{prompttable}
\prow{เติมชื่อ milestone ระดับแรกในช่อง 1.1-1.4 เรียงจากซ้ายไปขวาตามลำดับที่บรรลุ แล้วใส่ตัวอักษรงาน A-I ลงในช่องงานระดับถัดไปตามลำดับที่ต้องทำเพื่อให้บรรลุ milestone นั้น}
{Fill the first-level boxes 1.1-1.4 with milestone names, ordered from left to right as they are reached. Put task letters A-I in the second-level boxes, ordered by the work needed to reach each milestone.}
\prow{ใช้ทุกงานครั้งเดียว งานที่ทำขนานกันเรียงก่อนหลังในกลุ่มเดียวกันได้ ไม่เพิ่มงานหรือ milestone เพียงเพื่อเติมช่องให้เต็ม และเว้นช่องที่ไม่ใช้}
{Use every task exactly once. Parallel tasks may appear in either order within their group. Do not invent tasks or milestones merely to fill the template. Leave unused boxes empty.}
\end{prompttable}
\examwbs
\pair{พื้นที่ด้านล่างใช้ร่างหรืออธิบาย WBS ได้ คำตอบหลักต้องอยู่ในต้นไม้ข้างบน}
{You may use the space below for working or a brief explanation. Record the main answer in the tree above.}
\answerfill
\newpage
\continued{ข้อ 5 (ต่อ): บันทึกประชุมสำหรับ RACI}{Question 5 (continued): Meeting notes for RACI}
\bsub{5.2}{เติม RACI ตามบทบาทที่ได้รับมอบหมาย}{Complete the RACI chart}{2}
\pair{พิจารณาบทบาทของแต่ละคนในแต่ละงานจากบันทึกประชุมด้านล่าง กรอก R, A, C, I หรือ “-” เมื่อไม่มีบทบาทให้ครบทุกช่อง ไม่อนุมานบทบาทเพิ่มจากตำแหน่งงานเพียงอย่างเดียว}
{Use the meeting notes below to determine each person's part in each task. Enter R, A, C, I, or “-” for no part in every cell. Do not infer additional responsibilities solely from a person's job title.}
\begin{prompttable}
\prow{\textbf{E:} Aom ลงมือทำ PM รับผิดชอบผลสำเร็จ Bank ให้ความเห็นก่อนทำ Chai รับแจ้งหลังทำ Dao ไม่มีบทบาท}
{\textbf{E:} Aom performs the work; PM is accountable for the outcome; Bank is consulted beforehand; Chai is informed afterward; Dao has no part.}
\prow{\textbf{C:} Bank ลงมือทำ PM รับผิดชอบผลสำเร็จ Aom ให้ความเห็นก่อนทำ Chai รับแจ้งหลังทำ Dao ไม่มีบทบาท}
{\textbf{C:} Bank performs the work; PM is accountable for the outcome; Aom is consulted beforehand; Chai is informed afterward; Dao has no part.}
\prow{\textbf{G:} Chai ลงมือทำ PM รับผิดชอบผลสำเร็จ Bank ให้ความเห็นก่อนทำ Dao รับแจ้งหลังทำ Aom ไม่มีบทบาท}
{\textbf{G:} Chai performs the work; PM is accountable for the outcome; Bank is consulted beforehand; Dao is informed afterward; Aom has no part.}
\prow{\textbf{D:} Dao ลงมือทำ PM รับผิดชอบผลสำเร็จ Chai ให้ความเห็นก่อนทำ Aom รับแจ้งหลังทำ Bank ไม่มีบทบาท}
{\textbf{D:} Dao performs the work; PM is accountable for the outcome; Chai is consulted beforehand; Aom is informed afterward; Bank has no part.}
\end{prompttable}
\begingroup\renewcommand{\tabularxcolumn}[1]{m{#1}}
\begin{tabularx}{\linewidth}{|M{1.1cm}|*{5}{>{\centering\arraybackslash}X|}}\hline
Task & PM & Aom & Bank & Chai & Dao\\\hline
E &\rule{0pt}{1.25cm}&&&&\\\hline
C &\rule{0pt}{1.25cm}&&&&\\\hline
G &\rule{0pt}{1.25cm}&&&&\\\hline
D &\rule{0pt}{1.25cm}&&&&\\\hline
\end{tabularx}\endgroup
\newpage
\continued{ข้อ 5 (ต่อ): เงื่อนไขการจัดตาราง}{Question 5 (continued): Scheduling conditions}
\bsub{5.3}{วางกำหนดการเป็น Gantt chart}{Build the Gantt chart}{5}
\pair{\textbf{1 คะแนน (รวมอยู่ในข้อ 5.3):} อธิบายว่าคุณใส่ buffer ตรงใดและเพราะเหตุใด รวมทั้งเหตุผลที่ milestone ไม่ใช้เวลาทำงาน คำอธิบายต้องตรงกับตารางที่คุณวาด}
{\textbf{1 mark (included in Question 5.3):} Explain where you placed buffers and why, and why milestones take no working time. Your explanation must match your chart.}
\begin{prompttable}
\prow{โครงการเริ่มวันจันทร์ที่ 4 มกราคม 2027 ทำงานวันจันทร์ถึงศุกร์และไม่มีวันหยุดอื่น ใช้ owner และจำนวนวันของแต่ละงานตามตาราง A-I งานหนึ่งเริ่มได้ในวันทำงานถัดจากวันที่ predecessors และ buffers ของงานเหล่านั้นเสร็จครบ}
{The project starts on Monday 4 January 2027. Work is Monday to Friday with no other holidays. Use the owner and duration in the A-I task table. A task starts on the working day after all its predecessors and their buffers have finished.}
\prow{คุณและสมาชิกทีมทั้งสี่ทำ DormFix เต็มเวลา ส่วนผู้เกี่ยวข้องนอกทีมมีงานอื่นด้วย ใส่ task buffer 1 วันและ project buffer 2 วันเฉพาะจุดที่จำเป็น}
{You and all four team members work full-time on DormFix. External stakeholders also have other duties. Add a 1-day task buffer and a 2-day project buffer only where necessary.}
\prow{Milestone อยู่วันสุดท้ายที่ทำเครื่องหมายของงานสุดท้ายในกลุ่ม ส่วน milestone สุดท้ายอยู่วันสุดท้ายของ project buffer และไม่เพิ่มวันทำงานเพื่อทำ milestone}
{A milestone is reached on the last marked day of its group's final task. The final milestone is reached on the last day of the project buffer. A milestone does not add a working day.}
\prow{ใช้ตารางหน้าถัดไป กรอก No. ตาม WBS, Task, Owner, Start และ Due โดย Start/Due เป็นวันแรก/สุดท้ายที่ทำเครื่องหมาย รวม buffer ใช้ X แทนวันทำงาน, - แทน buffer, * แทน milestone และแสดง dependency arrows เว้นแถวที่ไม่ใช้}
{Use the chart on the next page. Fill No. from the WBS, Task, Owner, Start and Due. Start/Due are the first/last marked days, including buffers. Use X for work, - for a buffer, and * for a milestone. Draw dependency arrows and leave unused rows empty.}
\end{prompttable}
\answerlabel\answerfill
\newpage
% The only landscape sheet in the question booklet.
\begin{landscape}
\thispagestyle{empty}
\linespread{1}\fontsize{9.5}{12}\selectfont\setlength{\parskip}{2pt}
\noindent\makebox[\linewidth][l]{{\fontsize{9}{11}\selectfont ชื่อ-นามสกุล \hrulefill\quad รหัสนักศึกษา \makebox[1.9in]{\hrulefill}}}\par
\continued{ข้อ 5.3 (ต่อ): ตารางคำตอบ เดือนมกราคม 2027}{Question 5.3 (continued): January 2027 answer chart}
\pair{No. = รหัส WBS; Task = milestone หรือตัวอักษรงาน; Owner = ผู้ทำ (milestone ไม่มี owner); Start/Due รวม buffer; X = วันทำงาน, - = buffer, * = milestone; แสดงลูกศร dependency}
{No. = WBS identifier; Task = milestone or task letter; Owner = stated owner (none for a milestone); Start/Due include buffers; X = work, - = buffer, * = milestone. Draw dependency arrows.}
\examgantt
\vfill
\noindent{\fontsize{9}{11}\selectfont\textit{ข้อสอบรายวิชา CPE334 Software Engineering กลางภาคการศึกษาที่ \Semester}\hfill\textit{หน้าที่ \thepage/\pageref{LastPage}}}\par
\end{landscape}
\restoregeometry
\linespread{1.1}\fontsize{10.5}{14}\selectfont\setlength{\parskip}{4pt}
\bq{6}{Backlog grooming: จองสนามแบดมินตัน}{Backlog grooming: Book a badminton court}{10}
\begin{prompttable}
\prow{คุณทำหน้าที่ Product Owner ระหว่าง grooming ของ Issue \#12 “Book a badminton court” ฝ่ายกีฬาซึ่งเป็น client ต้องการให้นักศึกษาจองสนามครั้งละ 1 ชั่วโมง เพื่อไม่ต้องต่อคิวที่เคาน์เตอร์}
{You act as Product Owner while grooming Issue \#12, “Book a badminton court.” The sports office, your client, wants students to book a court for one hour so they do not have to queue at the counter.}
\prow{ระบบมีข้อมูลสนามและเวลาเปิดให้บริการแล้ว นักศึกษาแต่ละคนถือการจองได้ไม่เกิน 1 รายการต่อวัน ตัวอย่าง: สนาม 2 ว่างเวลา 17:00 วันนี้ หลังนักศึกษาจองสำเร็จ สนาม 2 เวลา 17:00 ต้องแสดง “Booked”}
{Courts and opening hours are already stored in the system. A student may hold only one booking per day. Example: Court 2 is free at 17:00 today. After a student books it, Court 2 at 17:00 must be shown as “Booked.”}
\prow{ทีมยังไม่ได้ตกลงผู้รับผิดชอบ milestone หรือ estimate ให้คงข้อมูลที่ได้รับ แยกข้อที่ต้องถามเพิ่มเติม และอย่าสมมติว่าทีมได้ตัดสินใจแล้ว}
{The team has not agreed an assignee, milestone or estimate. Preserve the given facts and distinguish follow-up questions from decisions that have actually been made.}
\end{prompttable}
\bsub{6.1}{เขียน user story}{Write the user story}{2}
\pair{เขียน role, function และ value ให้เชื่อมกันตามความต้องการ client}
{Write a user story connecting the role, requested function and value stated by the client.}
\blankbox{3cm}
\bsub{6.2}{ระบุ assumptions และ details}{Record assumptions and details}{2}
\pair{บันทึกสิ่งที่ทีมทราบแล้วจาก client แยกประเด็นที่ยังต้องถามออกจากข้อเท็จจริง และไม่เปลี่ยนความหมายของข้อจำกัดหนึ่งการจองต่อวัน}
{Record the facts already known from the client. Separate unresolved questions from facts, and retain the meaning of the one-booking-per-day limit.}
\answerfill
\newpage
\continued{ข้อ 6 (ต่อ): เกณฑ์ยอมรับและ metadata ของ Issue \#12}{Question 6 (continued): Acceptance criteria and Issue \#12 metadata}
\bsub{6.3}{เขียน Acceptance Criteria แบบ Gherkin}{Write Gherkin acceptance criteria}{4}
\pair{เขียน 2 สถานการณ์พร้อม Given, When และ Then: (1) จองสำเร็จตามตัวอย่างสนาม 2 เวลา 17:00 วันนี้ และ (2) ขอจองเพิ่มเมื่อมีการจองในวันนั้นอยู่แล้ว แต่ละสถานการณ์ต้องมีเงื่อนไขตั้งต้น การกระทำ และผลที่ตรวจได้สอดคล้องกัน}
{Write two scenarios using Given, When and Then: (1) successful booking of Court 2 at 17:00 today, and (2) requesting another booking while already holding a booking that day. Each scenario must connect its precondition, action and observable outcome.}
\pair{\textbf{สถานการณ์ที่ 1}}{\textbf{Scenario 1}}\blankbox{5cm}
\pair{\textbf{สถานการณ์ที่ 2}}{\textbf{Scenario 2}}\blankbox{5cm}
\bsub{6.4}{กำหนด metadata หลัง grooming}{Set the post-grooming metadata}{2}
\pair{กรอกทั้งสี่ช่องตามสิ่งที่ตกลงแล้ว เลือก label ให้สอดคล้องกับประเภทงาน และใส่ “-” ในช่องที่ยังไม่มีค่า}
{Complete all four fields using agreed information. Choose a label consistent with the type of work, and enter “-” where no value is available.}
\begin{tabularx}{\linewidth}{|P|P|P|P|}\hline
Assignees & Labels & Milestone & Estimate\\\hline
\rule{0pt}{1.7cm}&&&\\\hline
\end{tabularx}
\newpage
\bq{7}{Sprint planning: ShuttleNow}{Sprint planning: ShuttleNow}{5}
\begin{prompttable}
\prow{คุณเป็น Scrum Master ของ ShuttleNow แอปติดตามรถรับส่งในมหาวิทยาลัย Product Owner ระบุเป้าหมาย Sprint 6 ว่า “นักศึกษาติดตามรถรับส่งและทราบเวลาที่รถจะมาถึงได้”}
{You are Scrum Master of ShuttleNow, a campus shuttle tracking app. For Sprint 6, the Product Owner states the goal: “A student can track a shuttle and know when it arrives.”}
\prow{ทีมทำได้ 16 story points ใน Sprint ก่อน และใช้ velocity เดิม Backlog เรียงตามความสำคัญ Anchor คือ “ดูรายการเส้นทางรถรับส่ง” = 3 points ใช้ scale 1, 2, 3, 5, 8, 13, 21 การใหญ่ขึ้นหรือเล็กลงหนึ่งขั้นหมายถึงเลขถัดไปบน scale}
{The team completed 16 story points last sprint and retains that velocity. The backlog is ranked by priority. The anchor, “View the list of shuttle routes,” is 3 points. Use the scale 1, 2, 3, 5, 8, 13, 21; a step larger or smaller means the next number on that scale.}
\end{prompttable}
\bsub{7.1}{ประเมิน points และเลือกงาน}{Estimate points and select work}{3}
\pair{เติม points และ Yes/No ทุกแถว โดยใช้เป้าหมาย Sprint ข้อมูล capacity และลำดับความสำคัญที่ให้}
{Complete Points and Yes/No in every row, using the stated sprint goal, capacity and priority order.}
\begingroup\fontsize{9.5}{12.5}\selectfont
\renewcommand{\tabularxcolumn}[1]{m{#1}}
\begin{tabularx}{\linewidth}{|M{.55cm}|>{\raggedright\arraybackslash}X|M{1.35cm}|M{.8cm}|M{.9cm}|}\hline
No. & \bilingualcell{User story}{เรื่องผู้ใช้} & เทียบ anchor & Points & Yes/No\\\hline
1 & \bilingualcell{See a shuttle's live position}{ดูตำแหน่งรถสด} & +2 steps &\rule{0pt}{.9cm}&\\\hline
2 & \bilingualcell{Pay the fare by QR code}{จ่ายค่าโดยสาร QR} & Same &\rule{0pt}{.9cm}&\\\hline
3 & \bilingualcell{See the arrival time at my stop}{เวลาถึงป้ายฉัน} & +1 step &\rule{0pt}{.9cm}&\\\hline
4 & \bilingualcell{Receive an alert before arrival}{เตือนก่อนรถถึง} & Same &\rule{0pt}{.9cm}&\\\hline
5 & \bilingualcell{See arrival times for all route stops}{เวลาถึงทุกป้าย} & +1 step &\rule{0pt}{.9cm}&\\\hline
6 & \bilingualcell{Fix the route list being cut off on small phones}{แก้รายการเส้นทางถูกตัดบนมือถือ} & -1 step &\rule{0pt}{.9cm}&\\\hline
\end{tabularx}\endgroup
\bsub{7.2}{รวม points ของงานที่เลือก}{Total the selected story points}{1}
\blankbox{.7cm}
\newpage
\continued{ข้อ 7 (ต่อ): อธิบาย Sprint backlog ที่เลือก}{Question 7 (continued): Explain the selected sprint backlog}
\bsub{7.3}{อธิบายเหตุผลที่เลือกและไม่เลือก}{Explain your selection and exclusions}{1}
\pair{เชื่อมเหตุผลกับเป้าหมาย Sprint, velocity และลำดับงาน อธิบายให้ครบทั้งงานที่เลือกและไม่เลือก}
{Relate your reasoning to the sprint goal, velocity and order. Explain both selected and excluded work.}
\answerfill
\newpage
% Shared miniature specification. This exercise is self-contained.
\bpart{3}{(50 คะแนน) แผนภาพ การออกแบบหน้าจอ และ Git}{(50 marks) UML/UI and Git Workflow}
\continued{ข้อกำหนดร่วมสำหรับข้อ 8-10: ระบบแจ้งปัญหา TokTickIT}{Shared requirements for Questions 8-10: TokTickIT Service Tickets}
\pair{ใช้ข้อกำหนด R1-R12 ของ TokTickIT ด้านล่างตอบข้อ 8-10 โดยยึดเฉพาะกฎที่ให้ ไม่ต้องเพิ่มทุกฟังก์ชันจากแลปจริง}
{Answer Questions 8-10 using the self-contained TokTickIT specification R1-R12 below. Model only the stated behaviour; do not import other laboratory features.}
\begin{prompttable}
\prow{\textbf{R1. บัญชีและหน้าแรกตาม role} หลัง sign in ให้ Requester, IT Staff และ Administrator ไป Dashboard ตาม role; Admin ใช้ metrics ของ Staff และอาจเพิ่มจำนวนบัญชี Requester เห็นเฉพาะ Ticket ของตน User เป็น actor ทั่วไป; Requester/IT Staff สืบทอดจาก User และ Administrator สืบทอดจาก IT Staff Ticket Owner คือ IT Staff ที่ได้รับมอบหมาย ไม่ใช่ role แยก Server ต้องตรวจสิทธิ์แม้ซ่อน control}
{\textbf{R1. Accounts and role-based home pages.} After sign-in, route Requesters, IT Staff and Administrators to their role-specific Dashboards; Admin may add account counts to Staff metrics. Requesters see only their own Tickets. User is the general actor; Requester and IT Staff specialize User, and Administrator specializes IT Staff. Ticket Owner means assigned IT Staff, not a role. Enforce server-side authorization even when controls are hidden.}
\prow{\textbf{R2. สร้างและติดตาม Ticket} Requester สร้าง Ticket โดยกรอก title, description และ requested priority (LOW, MEDIUM หรือ HIGH) ต้องกรอกครบ ระบบสร้าง ticket number, เวลา และสถานะ NEW อัตโนมัติ เมื่อสร้างสำเร็จต้องเรียก Send Email ผ่าน Email Gateway ทุกครั้ง Requester ดู My Tickets และรายละเอียดได้เฉพาะของตน}
{\textbf{R2. Create and track a Ticket.} A Requester supplies the required title, description and requested priority (LOW, MEDIUM or HIGH). The system generates the ticket number, timestamps and NEW status. Every successful creation invokes Send Email through the Email Gateway. My Tickets and its detail view expose only that Requester's Tickets.}
\prow{\textbf{R3. ข้อมูลและผู้รับผิดชอบ Ticket} รายละเอียดมี ticket number, title, description, requester name/email, requested priority, IT priority, assigned staff (Ticket Owner), status, created time และ updated time ข้อมูลจากบัญชี เวลา หมายเลข และข้อความเดิมของ Requester อ่านได้อย่างเดียว Owner อาจยังไม่ถูกกำหนดหรือเป็น IT Staff ที่ยัง active ได้หนึ่งคน IT Staff เปลี่ยน IT priority, owner และสถานะได้ตามกฎของแต่ละรายการ}
{\textbf{R3. Ticket data and ownership.} Details show the ticket number, title, description, requester name/email, requested priority, IT priority, assigned staff (the Ticket Owner), status, created time and updated time. Account-derived data, timestamps, identifiers and the Requester's original text are read-only. A Ticket may be unassigned or have one active IT Staff member as its Owner. IT Staff may change IT priority, owner and status only as permitted for that Ticket.}
\prow{\textbf{R4. Dashboard และ Queue ของ IT Staff} Staff Dashboard สรุป Ticket ที่ไม่มี Owner, ที่ตนเป็น Owner, จำนวนตาม status/IT priority และรายการล่าสุด แต่ละ metric เปิด Queue ที่ตรงกัน Queue ค้น title และกรอง status/IT priority/Owner; แสดง ticket number, title, requester, status, requested/IT priority, Owner และ updated time เลือกแถวเปิด Form View ได้ ต้องมี pagination และแยก “ไม่มี Ticket” จาก “ไม่พบผลลัพธ์”}
{\textbf{R4. Staff Dashboard and Queue.} The Staff Dashboard shows unassigned and my-owned Tickets, counts by status/IT priority and recent updates; each metric opens a matching Queue. The Queue searches titles, filters by status/IT priority/Owner, and lists ticket number, title, requester, status, both priorities, Owner and updated time. Rows open Form View. Include pagination and distinguish no Tickets from no matches.}
\prow{\textbf{R5. การปรับ Ticket} Update Ticket ใช้เปลี่ยน IT priority หรือสถานะตามสิทธิ์ การย้าย owner เป็นพฤติกรรมเสริม Reassign Ticket ซึ่งเกิดเฉพาะเมื่อ IT Staff เลือกย้ายให้บุคลากรที่ยัง active การอัปเดตแบบอื่นทำได้โดยไม่ต้อง Reassign}
{\textbf{R5. Updates.} Update Ticket changes IT priority or an allowed status. Reassign Ticket is an optional extension performed only when IT Staff choose another active staff member as owner. Other updates do not require reassignment.}
\end{prompttable}
\newpage
\continued{ข้อกำหนดร่วมสำหรับข้อ 8-10 (ต่อ)}{Shared requirements for Questions 8-10 (continued)}
\begin{prompttable}
\prow{\textbf{R6. Action และไฟล์} IT Staff ใช้ Record Action บันทึกข้อความและผล OPEN หรือ COMPLETED โดยต้อง Send Email ทุกครั้งที่บันทึกสำเร็จ Attach File เป็นพฤติกรรมเสริมเฉพาะเมื่อเลือกไฟล์ และต้อง Validate Attachment เพื่อตรวจชนิด/ขนาดเสมอ บันทึก Action โดยไม่แนบไฟล์ได้}
{\textbf{R6. Actions and files.} Record Action stores work text and an OPEN or COMPLETED result, always invoking Send Email on success. Attach File optionally extends it when a file is selected, and always invokes Validate Attachment for type/size checks. An Action need not have a file.}
\prow{\textbf{R7. การสื่อสาร} Public Comments ให้ Requester ของ Ticket และ IT Staff อ่าน/เขียนได้ ส่วน Internal Notes ให้เฉพาะ IT Staff และ Administrator อ่าน/เขียน ผู้รับ Email ต้องมีสิทธิ์เห็นข้อความนั้นเสมอ Requester ไม่เห็น Internal Notes หรือจำนวนรายการภายใน}
{\textbf{R7. Communication.} A Ticket's Requester and IT Staff can read/write Public Comments. Only IT Staff and Administrators can read/write Internal Notes. Email recipients must be authorized to see its contents. A Requester sees neither Internal Notes nor their count.}
\prow{\textbf{R8. การแก้เสร็จ} Resolve Ticket ต้องเรียก Check Resolution Gate ทุกครั้ง เงื่อนไขคือมี owner มี Action ที่ COMPLETED อย่างน้อยหนึ่งรายการ และไม่มี Action ที่ยัง OPEN ถ้าขาดข้อใดให้ปฏิเสธ แสดงเหตุผล และคงสถานะเดิม}
{\textbf{R8. Resolution.} Resolve Ticket always invokes Check Resolution Gate. An owner, at least one COMPLETED Action, and no OPEN Actions are required. If any condition fails, refuse the operation, explain why, and retain the current status.}
\prow{\textbf{R9. เริ่มและพักงาน} Ticket ใหม่เป็น NEW; IT Staff claim แล้วเป็น OPEN; เริ่มทำงานจาก OPEN เป็น IN\_PROGRESS; ขอข้อมูลจาก IN\_PROGRESS เป็น WAITING\_FOR\_REQUESTER; เมื่อได้รับข้อมูลแล้ว IT Staff resume กลับเป็น IN\_PROGRESS การเขียน Comment อย่างเดียวไม่เปลี่ยนสถานะ}
{\textbf{R9. Start and pause work.} A created Ticket is NEW. Staff claim it to OPEN, start it to IN\_PROGRESS, request information to WAITING\_FOR\_REQUESTER, and resume it to IN\_PROGRESS after receiving the information. Posting a Comment alone does not change status.}
\prow{\textbf{R10. จบและยกเลิก} IT Staff resolve จาก IN\_PROGRESS เป็น RESOLVED เมื่อผ่าน R8; Requester ยอมรับผลเป็น CLOSED หรือ reopen เป็น IN\_PROGRESS ได้ Requester ยกเลิก Ticket ของตนจาก NEW หรือ OPEN เป็น CANCELLED เท่านั้น CLOSED และ CANCELLED เป็น terminal states ไม่อนุญาต transition อื่นนอกเหนือจากที่ระบุ}
{\textbf{R10. Finish or cancel.} IT Staff resolve IN\_PROGRESS to RESOLVED only after R8 passes. The Requester either accepts to CLOSED or reopens to IN\_PROGRESS. They can cancel their own NEW or OPEN Ticket to CANCELLED only. CLOSED and CANCELLED are terminal. No other transitions are permitted.}
\prow{\textbf{R11. Validation และ stale data} ข้อความว่างบันทึกไม่ได้ หาก version ที่ส่งเก่าให้คงค่าที่ผู้ใช้พิมพ์ แสดงสถานะล่าสุดและให้ Refresh โดยไม่ส่งซ้ำอัตโนมัติ แสดง field errors หรือเหตุผลที่บันทึกไม่ได้ใกล้ control ที่เกี่ยวข้อง}
{\textbf{R11. Validation and stale data.} Blank text cannot be saved. If the submitted version is stale, preserve the user's input, show the current record and offer Refresh without automatic resubmission. Show field errors or refusal reasons near the relevant control.}
\prow{\textbf{R12. ขอบเขตภาพ} ข้อ 8 เติม relationships ของ actors/use cases ที่ให้ รวม Sign In, หน้า Dashboard ตาม role และ inheritance ที่ระบุ ข้อ 9 เติม transitions ของ states ที่ให้ ข้อ 10 เติม Dashboard ของทั้งสาม role และหน้าจอ Queue/Detail ในกรอบ ไม่ต้องสร้างบัญชีใหม่ ระบบชำระเงิน หรือบริการภายนอกอื่น}
{\textbf{R12. Diagram scope.} In Question 8 add relationships among the supplied actors and use cases, including Sign In, role-specific Dashboards and the stated inheritance. In Question 9 connect the supplied states. In Question 10 complete the three role dashboards and the Queue/Detail screens in the provided frames. Do not add account registration, payment or other external services.}
\end{prompttable}
\newpage
\bq{8}{เติมความสัมพันธ์ใน Use Case Diagram}{Complete the Use Case Diagram relationships}{10}
\pair{\textbf{เกณฑ์คะแนน:} associations, routing ไป Dashboard ตาม role และขอบเขต 3; include 3; extend พร้อมเงื่อนไข 2; generalization และทิศทาง 2 รวม 10 คะแนน}
{\textbf{Marks:} associations, role-to-dashboard routing and scope 3; include 3; extend with conditions 2; generalization and direction 2. Total: 10 marks.}
\pair{ใช้ R1-R8 และ R12 เติมความสัมพันธ์ในโครงหน้าถัดไป Actors, use cases และ system boundary มีให้แล้ว ไม่ต้องวาด node เพิ่ม ให้ใช้ association เส้นทึบ, dependency เส้นประพร้อม \textless\textless include\textgreater\textgreater หรือ \textless\textless extend\textgreater\textgreater และ generalization เส้นทึบหัวสามเหลี่ยมโปร่งชี้ไป actor ทั่วไป}
{Using R1-R8 and R12, complete the skeleton on the next page. Actors, use cases and the system boundary are supplied; do not add nodes. Use solid associations, dashed dependencies labelled \textless\textless include\textgreater\textgreater or \textless\textless extend\textgreater\textgreater, and solid generalizations with hollow triangular arrowheads pointing to the more general actor.}
\begin{prompttable}
\prow{เชื่อม User กับ Sign In; เชื่อม Requester กับ Requester Dashboard, Create Ticket และ View My Tickets; เชื่อม IT Staff กับ Staff Dashboard, View Ticket Queue, Update Ticket, Record Action และ Resolve Ticket; เชื่อม Administrator กับ Admin Dashboard และ Maintain Accounts โดยไม่ต้องลากความสามารถที่สืบทอดซ้ำ เชื่อม Email Gateway กับ use case ที่ติดต่อบริการนี้}
{Connect User to Sign In; Requester to the Requester Dashboard, Create Ticket and View My Tickets; IT Staff to the Staff Dashboard, View Ticket Queue, Update Ticket, Record Action and Resolve Ticket; and Administrator to the Admin Dashboard and Maintain Accounts. Do not duplicate inherited associations. Connect Email Gateway to the use case that calls it.}
\prow{แยกพฤติกรรมที่ต้องเกิดทุกครั้งออกจากพฤติกรรมที่เกิดเฉพาะเมื่อเลือกแนบไฟล์หรือย้าย owner แสดงทิศหัวลูกศร stereotype และเงื่อนไขของ extension ให้ถูกต้อง แผนภาพนี้แสดงขอบเขตความสามารถ ไม่ใช่ลำดับงานหรือ transition ของ status}
{Distinguish mandatory reused behaviour from behaviour performed only when a file or reassignment is selected. Show arrow direction, stereotypes and extension conditions correctly. This diagram models capabilities and scope, not an execution sequence or status transitions.}
\end{prompttable}
\pair{\textbf{คำอธิบายหรือสมมติฐาน}}{\textbf{Explanation or assumptions}}\blankbox{4cm}
\newpage
\continued{ข้อ 8 (ต่อ): เติมเส้นและป้ายความสัมพันธ์ในโครงที่ให้}{Question 8 (continued): Add relationship edges and labels}
\examTicketUseCase
\newpage
\bq{9}{เติม State Diagram ของ Ticket}{Complete the Ticket State Diagram}{10}
\pair{ใช้ R8-R10 เชื่อม states ที่ให้ด้วยลูกศรแบบมีทิศทาง ระบุ event, ผู้กระทำ และ guard ที่จำเป็น แสดงจุดเริ่มและทางออกของ terminal states อย่าลากเส้น include, extend หรือ actor generalization ในภาพนี้ เพราะเป็น State Diagram}
{Using R8-R10, connect the supplied states with directed transitions. Label events, actors and required guards. Show the initial path and terminal-state exits. Do not use include, extend or actor generalization here: this is a State Diagram.}
\pair{Ticket Owner คือ IT Staff ที่ถูกกำหนดให้ Ticket นั้น ไม่ใช่ role ใหม่ ต้องแยก claim, start, request information, resume, resolve, accept, reopen และ cancel ให้ครบ claim กำหนด Owner; ใส่เงื่อนไข Resolve ของ R8 บนลูกศรหรือใช้ชื่อ guard แล้วอธิบายด้านล่าง การกระทำที่ไม่ผ่าน guard คงสถานะเดิม ไม่ต้องเพิ่มสถานะ Error}
{The Ticket Owner is the IT Staff member assigned to that Ticket, not a new role. Cover claim, start, request information, resume, resolve, accept, reopen and cancel. Claim assigns the Owner. Put R8's resolution conditions on the transition, or name the guard and define it below. A refused operation retains its state; do not add an Error state.}
\examTicketState
\pair{\textbf{นิยาม guard}}{\textbf{Guard definition}}\blankbox{1.2cm}
\newpage
\bq{10}{เติม Dashboard ตาม role และหน้าจอ Ticket}{Complete role dashboards and Ticket screens}{10}
\pair{ใช้ R1-R4 และ R7-R12 เติม Dashboard ของ Requester, IT Staff และ Administrator ในกรอบที่ให้ โดยหน้าแรกต้องตรงกับ role หลัง sign in ไม่ต้องวาดหน้า Login เติมข้อมูลลงกรอบ Queue/Detail ด้วย ใช้ * นำหน้า required fields, textbox สำหรับข้อมูลแก้ได้ และเส้นเฉียงสำหรับ read-only/calculated}
{Using R1-R4 and R7-R12, complete the Requester, IT Staff and Administrator Dashboards in the supplied frames. Each signed-in role must land on its own appropriate home; do not draw a Login screen. Complete the Queue/Detail frames as well. Mark required fields with *, editable values with textboxes, and read-only/calculated values with hatching.}
\ticketNotationReference
\bsub{10.1}{หน้าแรกและ Dashboard ตาม role}{Role-specific home dashboards}{3}
\pair{เติมองค์ประกอบและทางไปหน้ารายละเอียดให้ครบทั้งสามกรอบ Requester แสดง Ticket ของตนที่ยังเปิด รอข้อมูล เพิ่งอัปเดต และเพิ่งแก้เสร็จ; IT Staff แสดง Ticket ที่ยังไม่มี Owner, Ticket ที่ตนเป็น Owner, สรุปตาม status หรือ IT priority และรายการที่เพิ่งอัปเดต; Administrator ใช้สรุปของ Staff และอาจเพิ่มจำนวนบัญชี ทุก metric ที่กดได้ต้องพาไปยังรายการที่ตรงกัน}
{Complete all three frames with suitable content and drill-downs. The Requester view summarizes only their open, waiting, recently updated and recently resolved Tickets. IT Staff see unassigned Tickets, Tickets they own, counts by status or IT priority, and recently updated Tickets. The Administrator may reuse Staff metrics and add account counts. Each actionable metric must open the corresponding list.}
\examRoleDashboards
\newpage
\continued{ข้อ 10 (ต่อ): Ticket Queue ของ IT Staff}{Question 10 (continued): IT Staff Ticket Queue}
\bsub{10.2}{ค้นหาและกรอง Ticket ตาม Owner}{Search and filter Tickets by Owner}{3}
\pair{แสดงข้อมูล Queue ทุกช่องตาม R4 พร้อมค้นหา title, ตัวกรอง status/IT priority/Owner, การเลือกเปิด Ticket และ pagination แยก empty state จาก no-results และทำให้ล้างตัวกรองได้}
{Show every Queue field in R4, title search, status/IT priority/Owner filters, Ticket selection and pagination. Distinguish the empty state from no results and provide a way to clear filters.}
\examTicketQueue
\newpage
\continued{ข้อ 10 (ต่อ): Staff Ticket Detail}{Question 10 (continued): Staff Ticket Detail}
\bsub{10.3}{ข้อมูล เจ้าของ Ticket และการทำงาน}{Ticket data, Owner and actions}{4}
\pair{เติมข้อมูลตาม R3 แยก read-only จาก IT priority, Owner และ status ที่แก้ได้ Owner เลือกได้เฉพาะ IT Staff ที่ active แยก Public Comments กับ Internal Notes พร้อม Action และไฟล์แนบที่เลือกได้ แสดง feedback เมื่อ Resolve ไม่ผ่าน gate, ข้อความว่าง หรือ STALE\_UPDATE โดยคง input เดิมและไม่ส่งซ้ำอัตโนมัติ}
{Complete R3's data, separating read-only fields from editable IT priority, Owner and status. Owner selection must include only active IT Staff. Distinguish Public Comments from Internal Notes; provide Action and optional attachment entry. Show feedback for a blocked Resolve, blank text or STALE\_UPDATE, preserving unsaved input without automatic resubmission.}
\examTicketForm
\newpage
\continued{บริบทร่วมสำหรับข้อ 11-12: Git Workflow}{Shared context for Questions 11-12: Git Workflow}
\begin{prompttable}
\prow{ExamScore เป็นโปรแกรม Python แบบ command-line สำหรับอาจารย์ อ่านรายชื่อและคะแนนสอบจาก CSV ตรวจคะแนนช่วง 0-100 คำนวณสถิติและทำรายงานอัตราผ่าน เลือก Python เพื่อให้อ่านโค้ดและจัดการ CSV ได้โดยไม่รบกวนการฝึก Git}
{ExamScore is a small Python command-line program for an instructor. It reads names and scores from CSV, validates scores from 0 to 100, computes statistics and prepares a pass-rate report. Python keeps the code and CSV processing readable while the exercise focuses on Git.}
\prow{main แทน release ที่ deploy แล้ว sprint staging branches ใช้รวมและทดสอบงานที่เสร็จ ส่วน feature branches แยกการเปลี่ยนแปลงแต่ละงานจนพร้อมรวม}
{main represents the deployed release. Sprint staging branches collect completed work for integration and testing. Feature branches isolate individual changes until they are ready to merge.}
\end{prompttable}
\begin{tabularx}{\linewidth}{|>{\raggedright\arraybackslash}m{4.1cm}|P|}\hline
File & \bilingualcell{Purpose}{หน้าที่}\\\hline
\code{score\_data.py} & \bilingualcell{Loads and validates exam scores}{โหลดและตรวจสอบคะแนนสอบ}\\\hline
\code{scores.csv} & \bilingualcell{Sample student names and scores}{รายชื่อและคะแนนตัวอย่าง}\\\hline
\code{validation\_notes.txt} & \bilingualcell{New file documenting validation cases}{ไฟล์ใหม่บันทึกกรณีทดสอบ}\\\hline
\end{tabularx}
\begin{prompttable}
\prow{HEAD ระบุ local branch ที่ checkout ส่วน origin/* เป็น local remote-tracking reference ซึ่งบันทึกสถานะ GitHub ล่าสุดที่เครื่องนี้ทราบ การ switch ย้าย HEAD แต่ไม่ย้าย branch pointers}
{HEAD identifies the checked-out local branch. Each origin/* is a local remote-tracking reference recording the latest GitHub state known to this PC. Switching moves HEAD without moving branch pointers.}
\prow{สมมติ working directory สะอาด คำสั่งสำเร็จและไม่เกิด conflict เว้นแต่โจทย์ระบุเป็นอื่น ไม่มีการติดต่อ GitHub ถ้าไม่ได้แสดง fetch, push หรือ remote command แต่ละข้อเริ่มจาก graph ของตนเอง}
{Assume a clean working directory, successful commands and no conflict unless stated otherwise. No command contacts GitHub unless fetch, push or another remote command is shown. Each question starts from its own graph.}
\prow{\textbf{ใช้ชื่อไฟล์เมื่อ stage ห้ามใช้ \code{git add .}} แสดง local branches, origin/* และ HEAD ตามคำสั่งของแต่ละขั้น อย่านำ graph สุดท้ายของข้อ 11 ไปเริ่มข้อ 12}
{\textbf{Stage explicit file names; \code{git add .} is not permitted.} Show local branches, origin/* and HEAD as requested at each step. Do not use Question 11's final graph as Question 12's starting state.}
\prow{กราฟ A---B แทน B ที่มี parent A; ชื่อใน [ ] และ ref -> commit บอกตำแหน่งตัวชี้ HEAD -> branch คือ branch ปัจจุบัน เส้น / กับ \textbackslash{} แสดงการเชื่อม ancestry เครื่องหมาย > หน้าคำสั่งเป็น prompt คำตอบใช้กราฟตัวอักษรหรือภาพวาดได้}
{In A---B, B has parent A. Names in [ ] and ref -> commit locate references; HEAD -> branch marks the current branch. / and \textbackslash{} connect ancestry. A leading > is a command prompt. Answers may use a text graph or a drawn diagram.}
\end{prompttable}
\newpage
\bq{11}{แก้ defect แล้วรวมงานตรวจคะแนนเข้า sprint\_1}{Fix a defect and integrate score validation into sprint\_1}{10}
\pair{ExamScore เป็นเครื่องมือที่อาจารย์ใช้ตรวจคะแนนหลังสอบและจัดทำรายงานอัตราผ่าน รุ่นที่ deploy แล้วอยู่บน \code{main} ซึ่งชี้ไป commit C ส่วนทีม Sprint แยก \code{sprint\_1} จาก B และเพิ่ม D เพื่อปรับรูปแบบรายงาน ขณะเดียวกันทีมตรวจสอบคะแนนแยก \code{feature/load\_scores} จาก B และทำ parser refactor ไว้ที่ E ทั้งสองงานต้องอยู่ในประวัติเมื่อรวมรุ่นทดสอบ}
{ExamScore is a tool instructors use to validate marks and prepare a pass-rate report after an examination. The deployed release is on \code{main} at C. The Sprint team branched \code{sprint\_1} from B and added D for a report-format change. Separately, the score-validation team branched \code{feature/load\_scores} from B and committed a parser refactor as E. Both changes must remain in the test release.}
\pair{ก่อนปล่อยรุ่นทดสอบ ผู้ตรวจพบว่า \code{score\_data.py} ยังรับคะแนนต่ำกว่า 0 หรือสูงกว่า 100 ได้ ผู้พัฒนาแก้ validation บน feature branch แล้ว แต่การแก้ยังไม่ stage หรือ commit และ working tree มีการเปลี่ยนแปลงนี้เพียงไฟล์เดียว \code{scores.csv} ไม่ได้แก้ ต้องบันทึก patch นี้เป็น F แล้วรวม feature เข้า \code{sprint\_1} โดยเก็บงาน D ไว้ จากนั้นกลับไป \code{main} เพื่อตรวจรุ่นที่ deploy อยู่ ไม่มีการ push ในโจทย์นี้}
{Before preparing the test release, a reviewer finds that \code{score\_data.py} still accepts values below 0 or above 100. The developer has corrected validation on the feature branch, but the edit is not staged or committed; it is the only working-tree change. \code{scores.csv} is unchanged. Record the patch as F, integrate the feature into \code{sprint\_1} without losing D, then return to \code{main} to inspect the deployed release. Do not push.}
\begin{gitpanel}
A---B---C [main, origin/main]
     |\
     D E [HEAD -> feature/load_scores, origin/feature/load_scores]
sprint_1, origin/sprint_1 -> D
\end{gitpanel}
\pair{สมมติ working tree มีเฉพาะการแก้ \code{score\_data.py} ดังกล่าว ให้รันคำสั่งตามลำดับ ไม่มี conflict และ merge สร้าง commit เมื่อ branch แยกสายกัน}
{Assume the working tree contains only the stated edit to \code{score\_data.py}. Run the commands in order. There are no file conflicts, and a merge commit is created when the branches have diverged.}
\begin{gitpanel}
> git status
> git diff -- score_data.py
> git add -- score_data.py
> git commit -m "Reject out-of-range scores"
> git switch sprint_1
> git merge feature/load_scores
> git switch main
\end{gitpanel}
\bsub{11.1}{แสดง graph หลัง commit F}{Draw the graph after commit F}{2}
\pair{แสดง local branches, origin/* และ HEAD โดยไม่ย้าย refs ที่คำสั่งยังไม่เปลี่ยน}{Show all local branches, origin/* and HEAD; keep every reference that these commands have not moved.}
\blankbox{3.4cm}
\newpage
\continued{ข้อ 11 (ต่อ): รวม feature เข้า sprint\_1 และสลับกลับ main}{Question 11 (continued): Merge the feature into sprint\_1 and switch back to main}
\bsub{11.2}{แสดง graph หลัง merge เข้า sprint\_1}{Draw the graph after merging into sprint\_1}{5}
\pair{แสดง commit graph และ refs ทุกตัว ระบุชนิด merge, common ancestor และ parents ของ merge commit G อธิบายว่า local branch ใดขยับ และ origin/* ใดคงเดิม}
{Show the commit graph and every reference. Identify the merge type, common ancestor and both parents of merge commit G. Explain which local branch moves and which origin/* references stay unchanged.}
\blankbox{5.4cm}
\bsub{11.3}{แสดง graph หลัง switch ไป main}{Draw the graph after switching to main}{3}
\pair{ระบุ HEAD และตำแหน่ง branch/ref หลังคำสั่งสุดท้าย พร้อมอธิบายว่า switch สร้าง commit หรือไม่}
{State HEAD and every branch/reference after the final command. Explain whether switching creates a commit.}
\blankbox{3cm}
\newpage
\bq{12}{เขียนคำสั่งและวาด graph ทีละขั้น}{Write commands and redraw after each step}{10}
\pair{ข้อ 12 เริ่มใหม่ตามกราฟนี้: feature แยกจาก C และได้ทำ commit D กับ E ต่อกันไว้แล้ว ขณะนี้ HEAD อยู่ feature/load\_scores ทำ 12.1-12.3 ต่อเนื่องโดยคง refs ที่ไม่มีคำสั่งเปลี่ยน ไม่ใช้ผลจากข้อ 11}
{Question 12 starts afresh: the feature was branched at C and has since added D and E in sequence. HEAD is on feature/load\_scores. Complete 12.1-12.3 continuously, retaining unchanged refs. Do not carry over Question 11's result.}
\begin{gitpanel}
A---B [main, origin/main]
     \
      C [sprint_1_staging, origin/sprint_1_staging]
       \
        D---E [HEAD -> feature/load_scores]
\end{gitpanel}
\bsub{12.1}{เปลี่ยนไป sprint staging branch}{Switch to the sprint staging branch}{2}
\pair{เขียนคำสั่งเปลี่ยนจาก feature/load\_scores ไป sprint\_1\_staging แล้ววาด graph ใหม่ให้แสดงตำแหน่ง HEAD และ branch/remote-tracking references หลังคำสั่ง}
{Write the command to switch from feature/load\_scores to sprint\_1\_staging. Redraw the graph to show HEAD and all branch/remote-tracking references after the command.}
\pair{\textbf{คำสั่ง}}{\textbf{Command}}\blankbox{2cm}
\pair{\textbf{Graph หลังข้อ 12.1}}{\textbf{Graph after 12.1}}\answerfill
\newpage
\continued{ข้อ 12 (ต่อ): เริ่มจากผลของข้อ 12.1}{Question 12 (continued): Start from the result of 12.1}
\bsub{12.2}{รวม feature branch เข้า staging}{Merge the feature branch into staging}{3}
\pair{ขณะนี้อยู่ที่ sprint\_1\_staging ให้เขียนคำสั่งรวม feature/load\_scores เข้าสู่ branch ปัจจุบัน ระบุ merge type พร้อมเหตุผลจาก ancestry แล้ววาด graph หลัง merge โดยแสดง HEAD, local branches และ remote-tracking references ทุกตัว}
{While on sprint\_1\_staging, write the command to merge feature/load\_scores into the current branch. Identify the merge type and justify it from the ancestry. Redraw the graph after the merge, showing HEAD, all local branches and all remote-tracking references.}
\pair{\textbf{คำสั่งและประเภท merge}}{\textbf{Command and merge type}}\blankbox{4cm}
\pair{\textbf{Graph หลังข้อ 12.2}}{\textbf{Graph after 12.2}}\answerfill
\newpage
\continued{ข้อ 12 (ต่อ): เริ่มจากผลของข้อ 12.2}{Question 12 (continued): Start from the result of 12.2}
\bsub{12.3}{ตรวจการเปลี่ยนแปลง stage และ commit}{Inspect, stage and commit the file changes}{5}
\begin{prompttable}
\prow{หลังข้อ 12.2 ผู้พัฒนาใช้ editor แก้ \code{score\_data.py} และ \code{scores.csv} และสร้างไฟล์ใหม่ \code{validation\_notes.txt} การแก้เนื้อหาเกิดขึ้นแล้ว คุณไม่ต้องเขียนโค้ด Python หรือคำสั่ง editor}
{After 12.2, the developer uses an editor to modify \code{score\_data.py} and \code{scores.csv} and create \code{validation\_notes.txt}. The edits are already made; you do not need to write Python code or editor commands.}
\prow{เขียนคำสั่ง Git เพื่อตรวจการเปลี่ยนแปลง stage เฉพาะสามไฟล์นี้ และสร้าง commit F ด้วยข้อความ \code{Document score validation cases} จากนั้นวาด graph โดยแสดง F, HEAD และ references ทุกตัว ไม่มีคำสั่ง push ในขั้นนี้}
{Write Git commands to inspect the changes, stage only these three files, and create commit F with message \code{Document score validation cases}. Then draw the graph with F, HEAD and every reference. No push is performed in this step.}
\end{prompttable}
\pair{\textbf{คำสั่ง}}{\textbf{Commands}}\blankbox{5cm}
\pair{\textbf{Graph หลังสร้าง F}}{\textbf{Graph after creating F}}\answerfill
